\section{特征形式与乘性：Euler 乘积}
\label{sec:hecke-eigenforms}
% ============================================================================

\textbf{我们遇到了什么问题？}
前一节已经知道 $T_n\colon \Sk_k \to \Sk_k$ 是线性算子，$\{T_n\}$ 彼此交换。
但 $\Sk_k$ 的维数可以很大——以 $\SL_2(\Z)$ 为例，$\dim \Sk_{24} = 2$，
$\dim \Sk_{36} = 3$，以后级别升高维数会很快增大。
面对一个 $d$ 维空间和一族线性算子，我们能问的最自然的问题是：
\emph{有没有一组基，使得每个 $T_n$ 在这组基下都变成对角矩阵？}

线性代数的答案是：\textbf{如果一族算子彼此交换且都可对角化，
则它们有公共特征向量基}。每个公共特征向量 $f$ 满足
$T_n f = \lambda_n f$（$f$ 是 $T_n$ 的特征向量，特征值为 $\lambda_n$）。
在这样的基下，计算变得极其简单：
$T_n$ 的矩阵就是对角矩阵，对角元是各特征值 $\lambda_n$。

\textbf{真正的惊喜在哪里？}
假设 $f = \sum a_n q^n$ 是这样的特征向量，且 $a_1 = 1$（首项归一化）。
由 \cref{thm:hecke-on-q-exp} 的公式，比较 $T_n f$ 的第 1 个 Fourier 系数：
\[
  (T_n f)_1 = a_n.
\]
而 $T_n f = \lambda_n f$ 意味着 $(T_n f)_1 = \lambda_n \cdot a_1 = \lambda_n$。
所以 $\lambda_n = a_n$——\textbf{特征值 $\lambda_n$ 就等于 Fourier 系数 $a_n$！}

这个发现有深刻的算术推论：
由于 $T_m T_n = T_{mn}$（当 $\gcd(m,n)=1$），特征值的乘法 $\lambda_{mn} = \lambda_m \lambda_n$
等价于 Fourier 系数的乘法 $a_{mn} = a_m a_n$。
这就是 Ramanujan 的 $\tau(mn) = \tau(m)\tau(n)$（$\gcd(m,n)=1$）的\emph{概念上的}证明：
$\Delta$ 恰好是 $\Sk_{12}(\SL_2(\Z))$ 唯一的（归一化）特征形式，所以它的 Fourier 系数必须是乘性的。

更进一步，乘性的 Dirichlet 级数 $\sum a_n n^{-s}$ 天然分解成素数处局部因子之积（Euler 乘积），
这把 $f$ 与数论中最核心的对象——$L$-函数——联系起来了。
这正是整条费马大定理证明链的起点：模形式的 $L$-函数与椭圆曲线的 $L$-函数"相等"。

\begin{definition}[特征形式]
\label{def:eigenform}
设 $f \in \Sk_k$ 满足 $T_n f = \lambda_n f$ 对所有 $n \geq 1$。称 $f$ 为 \textbf{特征形式}
（eigenform 或 Hecke eigenform）。若进一步 $a_1(f) = 1$（首项归一化），
称 $f$ 为\textbf{归一化特征形式}（normalized eigenform）。
\end{definition}

\begin{theorem}[乘性定理]
\label{thm:multiplicativity}
设 $f = \sum a_n q^n$ 是归一化特征形式。则
\begin{enumerate}
  \item[(i)] \textbf{乘性}：若 $\gcd(m,n)=1$，则 $a_{mn} = a_m a_n$。
  \item[(ii)] \textbf{素数幂递推}：对素数 $p$ 和 $r \geq 1$，
        \[
          a_{p^{r+1}} = a_p \cdot a_{p^r} - p^{k-1} a_{p^{r-1}}.
        \]
\end{enumerate}
\end{theorem}

\begin{proof}
(i) 利用 Hecke 算子的代数关系：
\[
  T_m \circ T_n = T_{mn} \quad (\gcd(m,n)=1), \qquad
  T_{p^r} \circ T_p = T_{p^{r+1}} + p^{k-1} T_{p^{r-1}}.
\]
对归一化特征形式，$T_n f = a_n f$，故
\[
  a_{mn} f = T_{mn} f = T_m(T_n f) = T_m(a_n f) = a_n T_m f = a_n a_m f.
\]
两边除以 $f$（非零）得 $a_{mn} = a_m a_n$。

(ii) 由递推关系 $T_{p^{r+1}} = T_p T_{p^r} - p^{k-1} T_{p^{r-1}}$，作用在 $f$ 上：
\[
  a_{p^{r+1}} f = T_{p^{r+1}} f
  = T_p(a_{p^r} f) - p^{k-1}(a_{p^{r-1}} f)
  = (a_p a_{p^r} - p^{k-1} a_{p^{r-1}}) f.
\]
\end{proof}

\begin{example}[Ramanujan $\tau$ 函数的乘性]
\label{ex:tau-multiplicativity}
$\Delta(\tau) = q\prod_{n=1}^\infty(1-q^n)^{24} = \sum_{n=1}^\infty \tau(n)q^n$
是重量 12 的尖点形式，且 $\dim \Sk_{12}(\SL_2(\Z)) = 1$，
故 $\Delta$ 自动是（唯一的）归一化特征形式，从而：
\begin{itemize}
  \item $\tau(mn) = \tau(m)\tau(n)$，当 $\gcd(m,n)=1$；
  \item $\tau(p^{r+1}) = \tau(p)\tau(p^r) - p^{11}\tau(p^{r-1})$。
\end{itemize}
例如 $\tau(2)\tau(3) = (-24)(252) = -6048 = \tau(6)$ \checkmark。

Ramanujan 于 1916 年猜测 $\tau$ 是乘性函数；1937 年 Mordell 证明了这一点——
正是用 Hecke 算子（当时 Hecke 刚发展完该理论）。
\end{example}

\begin{theorem}[$L$-函数的 Euler 乘积]
\label{thm:L-function-euler-product}
设 $f = \sum_{n=1}^\infty a_n q^n$ 是重量 $k$、级别 $N$ 的归一化新特征形式。
定义
\[
  L(s, f) = \sum_{n=1}^{\infty} \frac{a_n}{n^s}, \quad \mathrm{Re}(s) \gg 0.
\]
则 $L(s,f)$ 有 Euler 乘积分解：
\[
  L(s, f)
  = \prod_{p \nmid N} \frac{1}{1 - a_p p^{-s} + p^{k-1-2s}}
    \prod_{p \mid N} \frac{1}{1 - a_p p^{-s}}.
\]
\end{theorem}

% TikZ figure: Euler product decomposition of L-function
% Shows how multiplicativity of a_n leads to Euler product L(f,s)
\begin{figure}[ht]
\centering
\begin{tikzpicture}[
  box/.style={draw, rounded corners=3pt, fill=blue!10, minimum width=2.6cm, minimum height=0.8cm, font=\small},
  arrow/.style={->, thick, >=stealth},
  label/.style={font=\small\itshape}
]

% Title node
\node[font=\bfseries\small] at (4, 3.6) {特征形式的 $L$-函数乘积展开};

% Left: Dirichlet series
\node[box, fill=orange!15] (dirichlet) at (0, 2) {$L(f,s) = \sum_{n=1}^{\infty} \frac{a_n}{n^s}$};

% Arrow
\draw[arrow] (dirichlet) -- node[above, label] {乘性} (3.5, 2);

% Right: Euler product
\node[box, fill=green!15, minimum width=3.6cm] (euler) at (5.8, 2) {$\prod_p \frac{1}{1 - a_p p^{-s} + p^{k-1-2s}}$};

% Below: two cases
\node[box, fill=blue!8, minimum width=2.8cm] (good) at (1.8, 0.3) {好素数 $p \nmid N$: $\frac{1}{1-a_p p^{-s}+p^{k-1-2s}}$};
\node[box, fill=red!8, minimum width=2.8cm] (bad) at (6.0, 0.3) {坏素数 $p \mid N$: $\frac{1}{1 - a_p p^{-s}}$};

\draw[arrow, dashed] (euler.south) -- ++(0,-0.4) -| (good.north);
\draw[arrow, dashed] (euler.south) -- ++(0,-0.4) -| (bad.north);

% Multiplicativity condition
\node[draw, dashed, rounded corners, fill=yellow!10, font=\small, text width=7cm, align=left]
  at (3.9, -1.5) {乘性：$\gcd(m,n)=1 \Rightarrow a_{mn}=a_m a_n$\\
  完全乘性：$a_{p^r} = a_p \cdot a_{p^{r-1}} - p^{k-1} a_{p^{r-2}}$};

\end{tikzpicture}
\caption{$L(f,s)$ 的 Euler 乘积分解：乘性的 Fourier 系数导致 $L$-函数写成素数处局部因子之积。}
\label{fig:eigenform-euler}
\end{figure}


\begin{proof}[证明思路]
\textbf{好素数处（$p \nmid N$）：}由乘性，把 Dirichlet 级数写成素数幂的乘积：
\[
  L(s,f) = \prod_p \sum_{r=0}^\infty \frac{a_{p^r}}{p^{rs}}.
\]
利用递推 $a_{p^{r+1}} = a_p a_{p^r} - p^{k-1}a_{p^{r-1}}$，
令 $X = p^{-s}$，对级数 $F_p(X) = \sum_{r=0}^\infty a_{p^r} X^r$ 做分析：
\begin{align}
  F_p(X) - a_p X F_p(X) + p^{k-1} X^2 F_p(X)
  &= a_0 + (a_1 - a_p a_0) X + \sum_{r=2}^\infty (a_{p^r} - a_p a_{p^{r-1}} + p^{k-1}a_{p^{r-2}}) X^r \\
  &= 1 + 0 \cdot X + 0 = 1.
\end{align}
故 $F_p(X) = 1/(1 - a_p X + p^{k-1} X^2)$。

\textbf{坏素数处（$p \mid N$）：}$a_{p^r} = a_p^r$（几何级数），故
$F_p(X) = 1/(1 - a_p X)$。
\end{proof}

\begin{remark}
Euler 乘积把 $L(s,f)$ 拆成无穷多个"\textbf{局部因子}"之积，
每个素数 $p$ 对应一个因子，编码了 $f$（或对应椭圆曲线）在 $p$ 处的局部算术信息。
这是现代数论的核心范式之一。在费马大定理的证明中，
谷山-志村猜想断言每条有理椭圆曲线的 $L$-函数等于某个模形式的 $L(s,f)$——
正是因为两者有相同的 Euler 乘积！
\end{remark}

\begin{example}[$\zeta$ 函数与重量 1]
当把上面的 Euler 乘积推广到更一般的 $L$-函数族时，
类比 Riemann $\zeta(s) = \prod_p (1-p^{-s})^{-1}$（"重量 0 的模形式"情形），
可以看到局部因子分母的次数刚好给出 $L$-函数的"\textbf{欧拉特征}"。
\end{example}

\subsection*{把全部 Hecke 算子打包：Hecke 代数}

到这里为止，我们一直一个一个地看 $T_n$。但从 Diamond 的观点看，更自然的对象其实是它们生成的整系数交换代数：单个算子给出单条递推，整座 Hecke 代数才真正组织起全部特征值、同余与局部化信息。

\begin{definition}[Hecke 代数]
\label{def:hecke-algebra}
对固定的权 $k$ 与级别 $N$，定义
\[
  \mathbb{T}_k(N)=\Z[T_1,T_2,T_3,\ldots]
  \subset \End_{\C}\bigl(\Sk_k(\GammaO(N))\bigr)
\]
为作用在 $\Sk_k(\GammaO(N))$ 上的\textbf{Hecke 代数}。由于 $T_1=\Id$ 且各 $T_n$ 两两交换，它是一个交换 $\Z$-代数。
\end{definition}

\begin{proposition}[有限秩性]
\label{prop:hecke-finite-rank}
$\mathbb{T}_k(N)$ 是有限秩、无挠的 $\Z$-模，而且
\[
  \rank_{\Z}\mathbb{T}_k(N)=\dim_{\C}\Sk_k(\GammaO(N)).
\]
\end{proposition}

\begin{proof}[证明思路]
空间 $\Sk_k(\GammaO(N))$ 的复维数有限，记为 $d$，因此 $\End_{\C}(\Sk_k(\GammaO(N)))$ 也是有限维复向量空间；于是由各 $T_n$ 生成的 $\mathbb{T}_k(N)$ 不可能有无限的 $\Z$-秩。另一方面，配对
\[
  \mathbb{T}_k(N)\times \Sk_k(\GammaO(N))\longrightarrow \C,
  \qquad (T,f)\longmapsto a_1(Tf)
\]
把 $T_n$ 送到“取第 $n$ 个 Fourier 系数”的泛函，因为 $a_1(T_nf)=a_n(f)$。借助 $q$-展开原理，这个配对足以把 Hecke 代数辨认成尖点形式空间积分结构的对偶，因此其秩恰好等于 $d$。
\end{proof}

\begin{definition}[局部化与完全化]
\label{def:hecke-completion}
设 $\mathfrak{m}\subset \mathbb{T}_k(N)$ 是极大理想。先在 $\mathfrak{m}$ 处局部化，再取 $\mathfrak{m}$-进完备化：
\[
  \mathbb{T}_{k,\mathfrak{m}}(N)
  = \varprojlim_r \bigl(\mathbb{T}_k(N)_{\mathfrak{m}}/\mathfrak{m}^r\bigr).
\]
它是一个\textbf{完备局部环}。把所有极大理想处的局部因子放在一起，得到完全化 Hecke 代数
\[
  \widehat{\mathbb{T}}_k(N)=\prod_{\mathfrak{m}} \mathbb{T}_{k,\mathfrak{m}}(N).
\]
\end{definition}

\begin{proposition}[局部因子的意义]
\label{prop:hecke-local-factors}
每个 $\mathbb{T}_{k,\mathfrak{m}}(N)$ 都记录了“模 $\mathfrak{m}$ 意义下彼此同余”的 Hecke 特征值系统。换言之，若两个特征形式的特征值在某个素数 $\ell$ 下同余，它们会落到同一个残差特征值系统，也就由同一个极大理想 $\mathfrak{m}$ 控制。
\end{proposition}

\begin{proof}[说明]
极大理想给出一个残差域 $\mathbb{T}_k(N)/\mathfrak{m}$，于是每个 $T_n$ 都有一个“模 $\mathfrak{m}$ 的特征值”。局部化把与这组残差特征值无关的部分都丢掉，完备化则进一步保留它的所有高阶同余信息。
\end{proof}

\textbf{为什么这对形变理论至关重要？} 给定一个模 $\ell$ 的残差 Galois 表示 $\bar\rho$，几乎所有好素数 $p\nmid N\ell$ 处的迹都应满足
\[
  \Tr\bigl(\bar\rho(\Frob_p)\bigr)\equiv T_p \pmod{\mathfrak{m}}.
\]
因此一个极大理想 $\mathfrak{m}$ 实际上选中了“一类模 $\ell$ 同余的模形式”。Wiles 在第 10 章里比较的对象，正是这样的局部 Hecke 代数 $\mathbb{T}_{k,\mathfrak{m}}(N)$ 与一个普形变环 $R_{\mathfrak{m}}$；著名的 $R=T$ 定理断言二者相同。也就是说，Hecke 代数不仅编码解析特征值，还同时编码了 Galois 表示的全部可允许形变。


% ============================================================================
